% ==STYLE-BEGIN==
\documentclass[UTF8,11pt]{ctexart}
\usepackage{amsmath,amssymb,amsthm}
\usepackage[a4paper,margin=2.4cm,top=2.3cm,bottom=2.3cm]{geometry}
\usepackage{fancyhdr}
\usepackage{enumitem}
\usepackage{hyperref}
\hypersetup{hidelinks}

% ---- 版式 ----
\linespread{1.15}
\setlength{\parskip}{0.35em}
\setlength{\parindent}{1.6em}
\setlist{nosep,leftmargin=2em}

% ---- 页眉 ----
\pagestyle{fancy}
\fancyhf{}
\fancyhead[L]{\small\coursename}
\fancyhead[R]{\small\thepage}
\renewcommand{\headrulewidth}{0.4pt}

% ---- 定理环境：中文用正体，标题加粗，按小节共用一条流水号 ----
\newtheoremstyle{cnthm}{0.7em}{0.7em}{\normalfont}{}{\bfseries}{}{0.5em}{}
\theoremstyle{cnthm}
\newtheorem{thm}{定理}[section]
\newtheorem{lem}[thm]{引理}
\newtheorem{prop}[thm]{命题}
\newtheorem{cor}[thm]{推论}
\newtheorem{defn}[thm]{定义}
\newtheorem{rem}[thm]{注}
\newtheorem{exa}[thm]{例}
\newtheorem{conj}[thm]{猜想}
\renewcommand{\proofname}{\bfseries 证明}
\providecommand{\coursename}{}
% ==STYLE-END==
\renewcommand{\coursename}{解析几何\,·\,1014}
\title{几何与代数 1014}
\xeCJKDeclareCharClass{CJK}{"2460 -> "2473} % 让 ①②③ 等圈号使用中文字体
\date{}
\begin{document}
\maketitle

\section{二次曲线的化简}

\[
a_{11}x^2+2a_{12}xy+a_{22}y^2+2b_1x+2b_2y+c = 0
\]
\noindent\textbf{1.} 利用坐标旋转
\[
\left \{ \begin{aligned}
\vec{e_1}' &= \cos{\theta} \vec{e_1} + \sin{\theta} \vec{e_2} \\
\vec{e_2}' &= -\sin{\theta} \vec{e_1} + \cos{\theta} \vec{e_2}
\end{aligned} \right. \\
\]
\[
\left \{ \begin{aligned}
x = \cos{\theta}x' - \sin{\theta}y' \\
y = \sin{\theta}x' + \cos{\theta}y'
\end{aligned} \right.
\tag{\text{①}}
\]
令
\[
\cot{2\theta} = \frac{a_{11}-a_{22}}{2a_{12}}
\]
则原式变成
\[
a_{11}'x'^2+a_{22}'y'^2+2b_1'x'+2b_2'y'+c' = 0
\]
\noindent\textbf{2.} 利用配方(坐标平移)将方程化为简单形式
\noindent\textbf{1.} 无混合项
\noindent\textbf{2.} 若有二次项,则无一次项
\noindent\textbf{3.} 若有一次项,则无常数项

\section{二次曲面的化简}

三元二次方程
\[
c_{11}x^2+c_{22}y^2+c_{33}z^2+2c_{12}xy+2c_{13}xz+2c_{23}yz+2b_1x+2b_2y+2b_3z+c = 0 \quad \tag{*}
\]
回顾空间内坐标变换公式
\[
\left \{ \begin{aligned}
x &= a_{11}x' + a_{12}y' + a_{13}z' \\
y &= a_{21}x' + a_{22}y' + a_{23}z' \\
z &= a_{31}x' + a_{32}y' + a_{33}z'
\end{aligned} \right.
\]
我们优先考虑消去二次项,平移量不会影响二次项,因此先不纳入考虑
将①带入$( * )$有
\[
c_{11}'x'^2+c_{22}'y'^2+c_{33}'z'^2+2c_{12}'x'y'+2c_{13}'x'z'+2c_{23}'y'z'+2b'_1x'+2b'_2y'+2b'_3z'+c'=0
\]
令
\[
F(x,y,z) = c_{11}x^2+c_{22}y^2+c_{33}z^2+2c_{12}xy+2c_{13}xz+2c_{23}yz
\]
代入①式有
\[
G(x',y',z') = c_{11}'x'^2+c_{22}'y'^2+c_{33}'z'^2+2c_{12}'x'y'+2c_{13}'x'z'+2c_{23}'y'z'
\]
令$f(x,y,z) = \frac{F(x,y,z)}{x^2+y^2+z^2}$,$g(x',y',z') = \frac{G(x,y,z)}{x'^2+y'^2+z'^2}$
此时$f(x,y,z) \rightarrow\text{①} g(x',y',z')$	
$f(x,y,z)$在$\mathbb{R}^3$值域与$f(x,y,z)$限制在单位球面$S^2 = \{ (x,y,z)|x^2+y^2+z^2=1 \}$的值域相同
于是,$f(x,y,z)$有最大值和最小值(空间上有界闭区间函数必有最大值/最小值)
不妨设$f(x,y,z)$在$P(x_0,y_0,z_0)\in S$取到最大值
取新的坐标系$[O;X',Y',Z']$s.t. z'轴是OP方向
此时,P点的新坐标为$(0,0,1)$,$g(x',y',z')$在$(0,0,1)$处取最大值
过$(0,0,1)$作平行于$x'$轴的直线$l$
$l$的方程是$\left \{ \begin{aligned} x' &= 0 +1 \cdot t \\ y' &= 0 + 0 \cdot t \\z' &= 1 + 0 \cdot 1 \end{aligned} \right.$,其中$t$是参数
令$\widetilde{g}(t) = g(x'(t),y'(t),z'(t))=\frac{c'_{11}t^2+2c'_{13}t+c'_{33}}{1+t^2}$

该函数在$t=0$处取最大值,故
\[
\tilde{g(0)}' = 0 \Rightarrow c'_{13} = 0
\]
同理,有$c_{23}' = 0$

由上,方程可化简为
\[
c_{11}'x'^2+c_{22}'y'^2+c_{33}'z'^2+2c_{12}'x'y'+2b'_{1}x'+2b'_{2}y'+2b'_{3}z'+c'
\]
再将坐标绕$z'$轴旋转,消去$x''y''$
于是,三元二次方程可化简成满足如下条件的形式
\noindent\textbf{1.} 无交叉项
\noindent\textbf{2.} 有平方项,则无一次项
\noindent\textbf{3.} 有一次项,则无常数项
\noindent\textbf{4.} 至多只有一个一次项
由上,我们可以得出二次曲面一共有17种

\end{document}

\end{document}
